\begin{table*}[htbp]
\centering
\scriptsize
\caption{Декомпозиция индекса стиля на стилевую и структурную части, две серии}
\label{tab:s14}
\begin{tabular}{>{\raggedright\arraybackslash}p{0.115\linewidth}>{\raggedright\arraybackslash}p{0.115\linewidth}>{\raggedright\arraybackslash}p{0.115\linewidth}>{\raggedright\arraybackslash}p{0.115\linewidth}>{\raggedright\arraybackslash}p{0.115\linewidth}>{\raggedright\arraybackslash}p{0.115\linewidth}>{\raggedright\arraybackslash}p{0.115\linewidth}>{\raggedright\arraybackslash}p{0.115\linewidth}}
\toprule
Контраст & Серия & Пар & Вклад common & Вклад format & Сумма & Δfull замороженный & Максимальное расхождение \\
\midrule
P3-P1 & v2 & 359 & +1.6740 & −3.4395 & −1.7655 & −1.7655 & 1.6e−04 \\
P2-P1 & v2 & 360 & −1.3927 & +4.6739 & +3.2812 & 3.2812 & 1.5e−04 \\
P3-P1 & v1 & 359 & +1.5458 & −3.4577 & −1.9118 & −1.9118 & 1.5e−04 \\
P2-P1 & v1 & 360 & −1.4879 & +4.6800 & +3.1921 & 3.1921 & 1.5e−04 \\
\bottomrule
\end{tabular}
\begin{minipage}{\linewidth}\footnotesize\vspace{2pt}
\textit{Примечание.} Единица анализа — пара текстов контраста. Знаменатель: 359 или 360 пар в зависимости от контраста. Данные: серии v2 и v1; v1 приводится как проверка устойчивости декомпозиции, а не как действующий результат. Разложение алгебраическое и точное: сумма вкладов совпадает с замороженным контрастом с точностью до 1.6·10⁻⁴, что и показывает колонка расхождения. Шкала: знак вклада читается по конвенции «больше значит более AI-подобно». Сокращения: common — стилевые категории признаков; format — структурная категория; Δfull — контраст полного балла из замороженного прогона. Пропуски: пропусков нет. Поправка на множественные сравнения не применялась.
\end{minipage}
\end{table*}
